%% kapa: exam=1 round=28 year=2017 period=1 form=A subject=real_estate subject_ko=부동산학원론 seq=33 number=113
%% source: Q-Net 감정평가사 기출문제 (cst00309) artlSeq=5199603 file=extracted/1008/5199603/2017년 제28회 감정평가사 1차시험 문제지/1교시 A형.pdf
%% conversion: kordoc 4.24.51 -> tools/gen1.py (auto); review=proofread-sonnet; flags=none
\begin{question}{113}
상업용 부동산 시장분석에 관한 설명으로 옳지 않은 것은?
\begin{choices}
  \item 소매점포 개설을 위한 시장분석의 절차는 부지평가 → 구역분석 → 시장선택의 3단계로 이루어진다.
  \item 통계적 분석방법은 기존통계를 분석해서 시장의 지역성을 포착하고, 그 지역성을 기초로 상권의 특성을 추계하는 방법이다.
  \item 상권추정기법에는 실제조사방법, 2차 자료 이용방법, 통계적 분석방법 등이 있다.
  \item 수정허프모델에서 고객의 구매확률은 상업지의 매장면적과 상업지로의 도달거리에 의해 결정된다.
  \item 체크리스트법은 매출액과 비용에 영향을 미칠 것으로 예상되는 다양한 요인들을 나열하고 이를 토대로 전문가적 경험에 의존하여 시장 내 대안부지들을 체계적으로 비교ㆍ평가하는 기법이다.
\end{choices}
\end{question}
\answer{1}
